\documentclass[../../main.tex]{subfiles}

\begin{document}

\section{Window-6 Countermodel Certificate and Equational Spectra}\label{app:window6-countermodel-certificate}

This appendix records the narrow certificate boundary for the window-\(6\)
finite countermodel used in the Equational Theories Project interface.  The
object is the affine-linear magma on \(\ZZ/21\ZZ\)
\[
    x\diamond y=7x+15y,
\]
or equivalently the same operation transported to the stable fiber \(X_6\) by
the stable-value ring equivalence.  The only upstream-facing certificate
claimed here is the compact finite entry
\[
    \texttt{Facts G [10, 52, 55] [43, 46]}.
\]
The surrounding Fibonacci and CRT structure is Lean-certified on the Omega
side; the ETP law satisfaction, spectrum count, and current database freshness
are verified computationally (see \texttt{\detokenize{scripts/equational_theory/}}).

\begin{lemma}[Stable-value affine transfer]\label{lem:window6-affine-stable-transfer}
Let \(m\geq 0\), and let \(X_m\) be the stable fiber equipped with
\(\mathrm{stableAdd}\) and \(\mathrm{stableMul}\).  Under the ring equivalence
\[
    X_m \simeq \ZZ/F_{m+2}\ZZ,
\]
every affine-linear magma operation
\[
    x\diamond y=ax+by
\]
on \(X_m\) is transported to the same coefficient formula over
\(\ZZ/F_{m+2}\ZZ\).  Hence a magma equation in this affine-linear family is
satisfied on \(X_m\) exactly when the transported coefficient congruences are
satisfied modulo \(F_{m+2}\).
\end{lemma}

\begin{proof}
The map \(\mathrm{stableValueRingEquiv}\) identifies \(X_m\) with
\(\ZZ/F_{m+2}\ZZ\), while \(\mathrm{stableAdd}\) and \(\mathrm{stableMul}\)
compute modular addition and multiplication on stable values.  For a variable
term the claim is immediate.  If \(t=u\diamond v\), then the transported value is
\[
    a\,\mathrm{val}(u)+b\,\mathrm{val}(v),
\]
which is exactly the recursive coefficient rule over the quotient ring.
Induction on the binary term proves commutation of evaluation with the
stable-value equivalence, and injectivity of the equivalence gives the stated
satisfaction criterion.  This is the formal Omega transfer bridge recorded in
\texttt{\detokenize{lean4/Omega/Folding/FiberRing.lean}} and
\texttt{\detokenize{lean4/Omega/Folding/FiberArithmetic.lean}}; it is not, by
itself, an upstream ETP theorem.
\end{proof}

\begin{proposition}[Window-\(6\) CRT rectangular band]\label{prop:window6-crt-rectangular-band}
For \(m=6\), the stable-value ring is
\[
    X_6\simeq \ZZ/21\ZZ\simeq \ZZ/3\ZZ\times \ZZ/7\ZZ.
\]
In this CRT presentation the operation \(x\diamond y=7x+15y\) becomes
\[
    \bigl((a_3,a_7),(b_3,b_7)\bigr)\longmapsto (a_3,b_7).
\]
Thus the window-\(6\) magma is the rectangular-band projection determined by
the complementary idempotents \(7\) and \(15\) in \(\ZZ/21\ZZ\).
\end{proposition}

\begin{proof}
The factorization \(F_8=21=3\cdot 7\) gives the CRT decomposition.  In
\(\ZZ/21\ZZ\), the Omega formalization proves
\[
    7^2=7,\qquad 15^2=15,\qquad 7\cdot 15=0,\qquad 7+15=1.
\]
Modulo \(3\), one has \(7\equiv 1\) and \(15\equiv 0\), so
\[
    7x+15y\equiv x\pmod 3.
\]
Modulo \(7\), one has \(7\equiv 0\) and \(15\equiv 1\), so
\[
    7x+15y\equiv y\pmod 7.
\]
Therefore the CRT image of \(x\diamond y\) keeps the \(\ZZ/3\ZZ\)-coordinate of
the left input and the \(\ZZ/7\ZZ\)-coordinate of the right input.  The
idempotent identities cited above are formalized in
\texttt{\detokenize{lean4/Omega/Folding/Window6.lean}}.
\end{proof}

\begin{lemma}[First--last term classifier]\label{lem:window6-first-last-classifier}
For every binary magma term \(t(x_1,\ldots,x_n)\), evaluation in the
window-\(6\) magma has the CRT form
\[
    t(x_1,\ldots,x_n)
    =
    \bigl(\pi_3(x_{\operatorname{first}(t)}),
          \pi_7(x_{\operatorname{last}(t)})\bigr),
\]
where \(\operatorname{first}(t)\) and \(\operatorname{last}(t)\) are the first
and last variables encountered in the standard left-to-right traversal of the
term.  Consequently a law \(s=t\) holds in this magma if and only if
\[
    \operatorname{first}(s)=\operatorname{first}(t)
    \quad\text{and}\quad
    \operatorname{last}(s)=\operatorname{last}(t).
\]
\end{lemma}

\begin{proof}
For a variable term, the first and last variables are both the variable itself,
and the displayed formula is tautological.  Suppose the formula is known for
\(u\) and \(v\).  By Proposition
\ref{prop:window6-crt-rectangular-band}, the value of \(u\diamond v\) is
obtained by taking the \(\ZZ/3\ZZ\)-coordinate from \(u\) and the
\(\ZZ/7\ZZ\)-coordinate from \(v\).  Hence the first variable of
\(u\diamond v\) is the first variable of \(u\), and the last variable of
\(u\diamond v\) is the last variable of \(v\).  This proves the formula by
structural induction.

If two terms have the same first and last variables, the formula gives equal
values for every assignment.  Conversely, if the first variables differ, choose
an assignment with distinct \(\ZZ/3\ZZ\)-coordinates on those variables; if the
last variables differ, choose an assignment with distinct
\(\ZZ/7\ZZ\)-coordinates.  The two CRT components have more than one element,
so the corresponding evaluations differ.
\end{proof}

\begin{theorem}[Window-\(6\) finite certificate]\label{thm:window6-fin21-facts-certificate}
Let \(G\) be the finite magma on \(\operatorname{Fin}(21)\) induced by
\[
    x\diamond y=7x+15y\pmod {21}.
\]
With respect to the ETP order-four equation list, \(G\) satisfies
\[
    E_{10},\qquad E_{52},\qquad E_{55},
\]
and refutes
\[
    E_{43},\qquad E_{46}.
\]
Equivalently, \(G\) provides the compact finite certificate
\[
    \texttt{Facts G [10, 52, 55] [43, 46]}.
\]
\end{theorem}

\begin{proof}
The algebraic form of the model is Proposition
\ref{prop:window6-crt-rectangular-band}, and the term classifier is Lemma
\ref{lem:window6-first-last-classifier}.  The correspondence between the ETP
indices \(10,43,46,52,55\) and the relevant term pairs is an external database
indexing fact, so the final indexed certificate is verified by direct
enumeration of the finite magma on \(\operatorname{Fin}(21)\).  This computation
checks all assignments for the listed laws and produces the stated
satisfied/refuted split; it is verified computationally (see
\texttt{\detokenize{scripts/equational_theory/}}).  No Automath/Omega Lean
theorem is claimed here for the ETP-indexed laws \(E_{10},E_{43},E_{46},E_{52}\),
or \(E_{55}\).
\end{proof}

\begin{corollary}[Equational spectrum and current-HEAD audit]\label{cor:window6-equational-spectrum}
For the same finite magma \(G\), the full ETP order-four spectrum contains
\[
    353
\]
satisfied equations and
\[
    4341
\]
refuted equations.  Hence \(G\) separates
\[
    353\cdot 4341=1,532,373
\]
ordered implication pairs.  At upstream ETP commit
\[
    \texttt{b99c0e486501e5b0690ef3fe5250d3aa57e63478},
\]
the combined certificate has no exact \(\texttt{Facts}\) entry and no
satisfied/refuted superset entry in \(\texttt{full\_entries.json}\).  Among the
six pairwise consequences of Theorem
\ref{thm:window6-fin21-facts-certificate}, the following five are also absent
as exact entries and as superset entries:
\[
\begin{array}{c|cc}
\text{pair} & \text{exact entries} & \text{superset entries}\\
\hline
E_{10}\not\Rightarrow E_{46} & 0 & 0\\
E_{52}\not\Rightarrow E_{43} & 0 & 0\\
E_{52}\not\Rightarrow E_{46} & 0 & 0\\
E_{55}\not\Rightarrow E_{43} & 0 & 0\\
E_{55}\not\Rightarrow E_{46} & 0 & 0
\end{array}
\]
The remaining pair \(E_{10}\not\Rightarrow E_{43}\) is already covered by two
\(\operatorname{Fin}(2)\) \(\texttt{SmallMagmas.lean}\) facts and is not part
of the nonredundancy claim.
\end{corollary}

\begin{proof}
The spectrum count is the direct truth-vector count over the \(4694\) indexed
ETP equations: \(353+4341=4694\), and every satisfied/refuted ordered pair is
an anti-implication witnessed by \(G\).  The current-HEAD novelty statement is
obtained by parsing the upstream \(\texttt{full\_entries.json}\) file, filtering
finite \(\texttt{Facts}\) entries, and checking both equality and
satisfied/refuted inclusion against the combined window-\(6\) certificate and
its pairwise projections.  The audit script copied with this paper is
\[
    \texttt{\detokenize{scripts/equational_theory/audit_window6_current.py}}.
\]
\end{proof}

\begin{remark}[Certificate boundary]\label{rem:window6-proof-boundary}
The Omega formalization certifies the stable-value ring equivalence, the
window-\(6\) CRT decomposition, and the idempotent arithmetic of \(7\) and
\(15\) in \(\ZZ/21\ZZ\).  The ETP-indexed law satisfaction, the \(353/4341\)
spectrum split, and the current upstream novelty audit are computational data
contributions.  They should be cited as a finite countermodel certificate, not
as a Lean proof of the ETP laws and not as a new graph-edge theorem beyond the
audited \(\texttt{Facts}\) contribution.
\end{remark}

\end{document}
